\documentclass[11pt]{amsart}

\usepackage[T1]{fontenc}
\usepackage{amsmath,amssymb,amsthm}
\usepackage[hidelinks]{hyperref}

\newtheorem{theorem}{Theorem}
\newcommand{\cl}{\operatorname{cl}}
\newcommand{\Int}{\operatorname{int}}
\newcommand{\bd}{\partial}

\title{A Short Proof of Bowron's ED-Space Conjecture}
\author{Chengyang Li}
\date{28 September 2026}

\begin{document}

\begin{abstract}
Bowron conjectured that every ED space, in the sense of his classification of Kuratowski operator monoids, contains a nonempty subset whose closure is the interior of its boundary. We give a short proof.
\end{abstract}

\maketitle

\section{The conjecture}

We follow the operator notation of Bowron~\cite{Bowron2024}: for a subset $A$ of a topological space,
$bA=\cl A$, $iA=\Int A$, and $fA=\bd A$.
Composition is written by juxtaposition, so for example $ibA=\Int(\cl A)$.

In the Closing Remarks of~\cite{Bowron2024}, Bowron conjectured that every ED space contains a nonempty subset $A$ such that
\[
 bA=ifA.
\]
Here ``ED space'' is used in the sense of Bowron's classification. In particular, the corresponding global operator identities include
\[
 bib=ib,\qquad bib\ne bi.
\]

\begin{theorem}
Every ED space in the sense of~\cite{Bowron2024} contains a nonempty subset $A$ such that
\[
 bA=ifA.
\]
\end{theorem}

\begin{proof}
Since $bi\le bib=ib$ as operators and $bib\ne bi$, there is a subset $B$ such that
\[
 biB\subsetneq ibB.
\]
By Lemma 8(i) of~\cite{Bowron2024},
\[
 ifB=ibB\setminus biB.
\]
Hence
\[
 U:=ifB
\]
is nonempty.

The set $U$ is open by definition. Since $fB$ is closed, applying the identity $bib=ib$ to $fB$ gives
\[
 b(ifB)=ifB,
\]
so $U$ is closed as well. Thus $U$ is clopen.

Let
\[
 A:=B\cap U.
\]
Since $A\subseteq U$ and $U$ is closed, $bA\subseteq U$.
Conversely, let $x\in U$ and let $V$ be an open neighbourhood of $x$. Then $V\cap U$ is an open neighbourhood of $x$. Moreover,
\[
 U\subseteq fB\subseteq bB,
\]
so $(V\cap U)\cap B\ne\varnothing$. Hence $V\cap A\ne\varnothing$.
Thus every neighbourhood of $x$ meets $A$, and therefore $x\in bA$. We conclude that
\[
 bA=U.
\]

We next claim that $iA=\varnothing$. Otherwise some nonempty open set $O$ would satisfy $O\subseteq A$. Since $O\subseteq B$, this would imply $O\subseteq iB$; but also $O\subseteq U\subseteq fB$, contradicting
\[
 iB\cap fB=\varnothing.
\]
Therefore $iA=\varnothing$, and hence
\[
 fA=bA\setminus iA=U.
\]
Finally, $U$ is open, so
\[
 ifA=iU=U=bA.
\]
Since $U\ne\varnothing$ and $bA=U$, the set $A$ is nonempty.
\end{proof}

\section*{AI assistance}
The initial argument was obtained with assistance from OpenAI's GPT. The author subsequently checked and reformulated the proof and takes responsibility for the mathematical content of this note.

\section*{Acknowledgement}
The author thanks Mark Bowron for helpful correspondence.

\begin{thebibliography}{1}

\bibitem{Bowron2024}
M. Bowron,
\emph{Boundary-border extensions of the Kuratowski monoid},
Topology Appl. \textbf{341} (2024), 108703.
\href{https://doi.org/10.1016/j.topol.2023.108703}{doi:10.1016/j.topol.2023.108703}.
Preprint: \href{https://arxiv.org/abs/2210.10928}{arXiv:2210.10928}.

\end{thebibliography}

\end{document}
